{\kaishu\zihao{4} 参照以下提纲撰写，要求内容翔实、清晰，层次分明，标题
  突出。\cemph{请勿删除或改动下述提纲标题及括号中的文字。}}

\part{立项依据与研究内容（建议8000字以下）：}
\label{part:background}

\section{项目的立项依据}{（研究意义、国内外研究现状及发展动态分析，需
  结合科学研究发展趋势来论述科学意义；或结合国民经济和社会发展中迫切需
  要解决的关键科技问题来论述其应用前景。附主要参考文献目录）；}
\label{sec:background}

\subsection{研究意义}
\label{sec:motivation}

略。

\subsection{国内外研究现状}
\label{sec:review}

可以引用相关文献~\cite{bengio_representation_2013}对研究意义与国内外研究现状进行说明。

\printbibliography[segment=\therefsegment,heading=reftype]

\section{项目的研究内容、研究目标，以及拟解决的关键科学问题}{（此部分为重点阐述内容）；}
\label{sec:proposals}

\subsection{研究内容}

略。

\begin{figure}[h]
  \centering

  \caption{测试图表标题。}
  \label{fig:test}
\end{figure}

\begin{table}[h]
  \centering
  \begin{tabular}{cc}
    \hline
    标题 & 内容 \\
    \hline
    科目1 & 内容1 \\
    \hline
  \end{tabular}
  \caption{测试图表标题。}
  \label{tab:test}
\end{table}

\subsection{研究目标}
\label{sec:goals}

略。

\subsection{拟解决的关键科学问题}

略。

\section{拟采取的研究方案及可行性分析}{（包括研究方法、技术路线、实验手段、关键技术等说明）；}
\label{sec:tech-feasibility}

略。

\section{本项目的特色与创新之处；}{}
\label{sec:innovations}

略。

\section{年度研究计划及预期研究结果}{（包括拟组织的重要学术交流活动、国际合作与交流计划等）。}
\label{sec:plans}

略。


\part{研究基础与工作条件}
\label{part:foundations}

\section{研究基础}{（与本项目相关的研究工作积累和已取得的研究工作成绩）；}
\label{sec:foundatioins}

\newrefsegment

  申请人针对\textbf{某问题}进行了研究~\cite{xia_saliency_2015}。

  在使用\texttt{biblatex} 3.4以上版本的情况下，建议使用
  \texttt{biblatex}提供数据标注功能。如第二个作者是自己，用粗体标记自己的姓
  名；且为文章的通记作者，加上角标``*''进行标注。则仅需要在bibtex文
  件的相应条目中增加以下标记：
\begin{verbatim}
  author = "Xia, C and Qi, F and Shi, G and Wang, P",
  author+an = {2=self;2:family=corr},
\end{verbatim}

  需要强调的作者列表有三种形式，分别适用于不同版本的\texttt{biblatex}。
  \begin{description}
  \item[3.4+ (20160301)] 请使用\texttt{biblatex}中提供的数据标注(Data
    annotation)功能。
  \item[3.3 (20160301)] 由于姓名处理的宏更新，请使用哈希字符串列表指定需要强调
    的作者。
\begin{verbatim}
    \initauthors{{72b3cccfc646adeb1d6b20320b56fd7d}}
\end{verbatim}
  \item[3.2- (20151228)] 使用旧式的姓名处理宏，使用下列形式的作者列表。
\begin{verbatim}
    \forcsvlist{\listadd\boldnames}{{Qi, F\bibinitperiod}}
\end{verbatim}
  \end{description}
  其中命令\texttt{initauthors}中的所需要的姓名的哈希字符串可以
  在Biber/bibtex 生成的文件\texttt{*.bbl} 中找到。

\printbibliography[segment=\therefsegment,heading=cvtype,title={相关工作}]

\section{工作条件}{（包括已具备的实验条件，尚缺少的实验条件和拟解决的
  途径，包括利用国家实验室、国家重点实验室和部门重点实验室等研究基地的
  计划与落实情况）；}
\label{sec:condition}

略。

\section{正在承担的与本项目相关的科研项目情况}{（申请人和项目组主要参与
  者正在承担的与本项目相关的科研项目情况，包括国家自然科学基金的项目和
  国家其他科技计划项目，要注明项目的名称和编号、经费来源、起止年月、与
  本项目的关系及负责的内容等）；}
\label{sec:projects}

略。

\section{完成国家自然科学基金项目情况}{（对申请人负责的前一个已结题科学
  基金项目（项目名称及批准号）完成情况、后续研究进展及与本申请项目的关
  系加以详细说明。另附该已结题项目研究工作总结摘要（限500字）和相关成果
  的详细目录）。}
\label{sec:finished-project}

\newrefsegment

略。

\part{其他需要说明的问题}
\label{cha:others}

\section{}{申请人同年申请不同类型的国家自然科学基金项目情况（列明同年
  申请的其他项目的项目类型、项目名称信息，并说明与本项目之间的区别与联
  系）。}

无

\section{}{具有高级专业技术职务（职称）的申请人或者主要参与者是否存在
  同年申请或者参与申请国家自然科学基金项目的单位不一致的情况；如存在上
  述情况，列明所涉及人员的姓名，申请或参与申请的其他项目的项目类型、项
  目名称、单位名称、上述人员在该项目中是申请人还是参与者，并说明单位不
  一致原因。}

无

\section{}{具有高级专业技术职务（职称）的申请人或者主要参与者是否存在与
  正在承担的国家自然科学基金项目的单位不一致的情况；如存在上述情况，列
  明所涉及人员的姓名，正在承担项目的批准号、项目类型、项目名称、单位名
  称、起止年月， 并说明单位不一致原因。}

无

\section{}{其他。}

无

使用方法参见样例文件 \texttt{my-proposal.tex}。

